
\chapter{Подписки и Card-on-File}~\label{ch:subscriptions}
Подписка начинается с первого согласия на хранение реквизитов. Для
продукта январский платёж и февральское
автосписание выглядят как продолжение одной услуги, но для эмитента
это две разные авторизационные ситуации. Если связь между ними
разорвать, эмитент теряет контекст операции, схема начинает считать
её подозрительной, а клиент получает лишние отказы там, где платёж
должен был пройти спокойно.

\section{Stored Credential Framework}~\label{sec:subscriptions-scf}

До появления единых правил мерчанты решали повторные списания каждый
по-своему: кто-то хранил PAN, кто-то полагался на токенизацию, кто-то
каждый месяц отправлял обычный запрос и надеялся, что эмитент сам
угадает контекст. Эмитенту не сообщали главное: это продолжение уже согласованной
цепочки или новая операция без явного участия клиента.

Отсутствие контекста делало \enquote{тихую} MIT неотличимой от обычного
CNP-запроса: мерчант мог начать регулярные списания без явного
согласия клиента или продолжить их после его отмены, и эмитент
не имел технической возможности это различить. Именно этот разрыв
и закрыл SCF: контекст первого согласия стал обязательным полем
авторизационного сообщения, а не опциональным пожеланием.

Visa и Mastercard сформировали фреймворк сохранённых реквизитов
(Stored Credential Framework, SCF), набор правил, по которым схема
различает первое согласие клиента и последующие списания так, чтобы
контекст первого согласия стал видимым в авторизационном
сообщении~\cite{visa-scf,mc-scf}. Карта остаётся у мерчанта как сохранённый реквизит, но
каждая операция должна ясно показывать, каким образом она связана с
исходным согласием клиента.

В авторизационном сообщении эта связь выражается несколькими
признаками: первичным или последующим использованием сохранённых
реквизитов, типом сценария и ссылкой на исходную авторизацию, связывающей
всю цепочку~\cite{mc-tpr,cybs-initial-transactions}.
SCF делает видимой границу между
первичным согласием клиента и дальнейшим автосписанием.

На уровне авторизационного сообщения SCF сводится к нескольким
конкретным индикаторам. Visa использует Stored Credential
Transaction Indicator: значение \texttt{initial} при первой
CIT и \texttt{subsequent} при последующей MIT, а рядом передаётся
тип инициатора (cardholder-initiated или merchant-initiated)
и причина MIT (recurring, installment,
unscheduled)~\cite{visa-scf,cybs-initial-transactions}. Mastercard
кодирует ту же информацию через Payment Initiation Channel
в DE~48 и требует передавать Trace~ID (DE~48 subelement~63)
исходной CIT в каждой последующей MIT~\cite{mc-scf,mc-tpr}.

Пример: клиент оформляет месячную подписку. Первый платёж идёт
как CIT с indicator = \texttt{initial} и
получает от схемы идентификатор транзакции (NTID у Visa, Trace~ID
у Mastercard). Через 30 дней мерчант отправляет MIT с
indicator = \texttt{subsequent}, типом \texttt{recurring} и
ссылкой на NTID первой CIT. Если мерчант не передаст ссылку
на исходную авторизацию или укажет неверный тип, эмитент увидит
запрос без контекста и вправе применить стандартную CNP-политику --
вплоть до отказа. Схема при аудите может квалифицировать такую
операцию как нарушение правил SCF~\cite{mc-tpr,visa-core-rules}.

\begin{figure}[tbp]
  \centering
  \includefiguresvg[width=.8\linewidth]{assets/figures/ch22-subscriptions/scf-lifecycle}
  \caption{Жизненный цикл SCF: CIT создаёт согласие, MIT продолжает цепочку.}~\label{fig:scf-lifecycle}
\end{figure}

Ограничение этого подхода тоже существенно: SCF не делает любую повторную
операцию автоматически безопасной. Если мерчант не передаёт контекст
или путает тип операции, эмитент вправе трактовать запрос как обычную
бесконтекстную CNP-транзакцию и применить к нему более строгую политику
проверки.

\section{CIT и MIT}~\label{sec:subscriptions-cit-mit}

Подписочный процесс почти всегда состоит из двух сцен. В первой клиент
сам оформляет услугу, соглашается на сохранение реквизитов и проходит
все необходимые шаги подтверждения. Во второй мерчант через неделю или
месяц пытается списать очередной платёж уже без участия клиента. Именно
на границе между этими двумя сценами чаще всего и ломается подписочный
процесс.

Транзакция, инициированная клиентом (customer-initiated transaction,
CIT), относится к первой сцене. Клиент сам нажимает кнопку оплаты,
вводит данные карты, проходит аутентификацию, если она требуется
правилами рынка. Так выглядит первый платёж за подписку, оформление
Card-on-File или покупка, после которой мерчант получает право
сохранить реквизиты. В этот момент клиент платит, и система фиксирует
исходную точку цепочки, получая связующий идентификатор исходной
авторизации (в документации
  сети или шлюза он называется transaction identifier, network transaction
ID, \texttt{previousTransactionID} или Trace ID)~\cite{mc-tpr,cybs-initial-transactions,cybs-mc-subscription-mit}.

Согласие клиента возникает в CIT, но сохранение реквизитов само по
себе ещё не означает право на любые будущие списания. Дальше система
должна хранить этот идентификатор и передавать его в последующих
операциях, иначе связь между первым согласием и следующими списаниями
быстро исчезнет из авторизационного контекста~\cite{mc-tpr,cybs-initial-transactions}.

Транзакция, инициированная мерчантом (merchant-initiated transaction,
MIT), относится уже ко второй сцене. Это списание без активного участия
клиента: ежемесячная подписка, доначисление по уже оказанной услуге или
платёж по ранее согласованной схеме. Здесь мерчант инициирует запрос,
но обязан показать эмитенту, что операция продолжает уже существующую
цепочку. Поэтому MIT почти всегда живёт рядом со SCF, схемными
признаками типа операции и передачей идентификатора исходной
авторизации~\cite{mc-tpr,cybs-mc-subscription-mit}.

Ограничение здесь жёсткое: MIT без правильного контекста может быть
расценена как новая и подозрительная транзакция. В европейском трафике
это особенно заметно, потому что при удалённом установлении MIT-мандата
сильная аутентификация клиента (strong customer authentication, SCA)
применяется на этапе его создания, а
последующие MIT должны быть корректно помечены и связаны с исходной
авторизацией~\cite{eba-mit-qa,mc-tpr}. Если инициатор и
контекст описаны неверно, эмитент вправе отклонить платёж даже тогда,
когда бизнес-логика мерчанта считает его штатным.

CIT создаёт согласие, а MIT использует уже созданную цепочку.
Для подписочного бизнеса правило простое: первый платёж должен пройти
и заложить техническую основу для всех следующих списаний.

\section{Идентификатор исходной авторизации и цепочка согласия}~\label{sec:subscriptions-ntid}

Токен отвечает на вопрос, каким реквизитом платить, а идентификатор
исходной авторизации отвечает на вопрос, с каким ранее данным
согласием связано текущее автосписание. Это разные
сущности, и смешивать их опасно: токен можно обновить при перевыпуске
карты, а идентификатор цепочки должен храниться отдельно от самих
реквизитов. Mastercard в ЕЭЗ, Великобритании и Гибралтаре прямо требует
в последующих запросах авторизации по регулярным платежам передавать
уникальный Trace ID из ответа на исходную авторизацию, а эмитент должен
уметь по нему восстановить исходную транзакцию~\cite{mc-tpr}. В
практической документации Visa/Cybersource та же идея сформулирована
ещё проще: после начальной операции нужно сохранить идентификатор
транзакции (transaction identifier), потому что он понадобится для
последующих операций~\cite{cybs-initial-transactions}.

Терминология различается между схемами. Visa называет
этот идентификатор Network Transaction ID (NTID) и
возвращает его в ответе на авторизацию; Mastercard использует
термин Trace~ID и требует его передачу в DE~48 последующих
авторизаций~\cite{mc-tpr,cybs-initial-transactions}. У
Cybersource и других шлюзов поле может называться
\texttt{previousTransactionID} или
\texttt{networkTransactionID} -- разные имена для одной
сущности.

Хранить идентификатор нужно на уровне подписки, а не на уровне
транзакции. Каждый последующий MIT ссылается на одну и ту же
исходную CIT, поэтому идентификатор живёт столько же, сколько
сама подписка. Если система хранит его только в записи
последней транзакции, при любом сбое записи цепочка обрывается.
На практике это означает отдельное поле в модели подписки,
которое заполняется при первой CIT и не перезаписывается
последующими ответами~\cite{cybs-mc-subscription-mit}.

Типичный сбой возникает при миграции между PSP. Мерчант переносит
токены и историю подписок, но забывает перенести идентификатор
исходной авторизации первой CIT. На следующем цикле MIT приходит уже
без этого якоря, и мерчант теряет связующий идентификатор
между первым согласием и последующим списанием~\cite{mc-tpr,cybs-mc-subscription-mit}.

\begin{figure}[tbp]
  \centering
  \includefiguresvg[width=.8\linewidth]{assets/figures/ch22-subscriptions/ntid-chain}
  \caption{Цепочка NTID: от первой CIT через MIT к точке обрыва при миграции PSP.}~\label{fig:ntid-chain}
\end{figure}

На бумаге защита выглядит просто: нужно хранить этот идентификатор как
обязательный атрибут подписки, передавать его во всех последующих MIT и
не разрывать цепочку при смене процессора. На практике это требует
такой же аккуратности, как работа с токенами и самой записью о
подписке. Если исходная CIT не была оформлена корректно, никакой
последующий косметический ремонт уже не спасёт цепочку согласия.

Идентификатор исходной авторизации -- это ссылка на согласие, из
которого вырастает вся последующая цепочка.

\section{Мягкий отказ и повторная попытка}~\label{sec:subscriptions-soft-decline}

Когда MIT получает отказ, который не закрывает сценарий окончательно,
важна реакция: схема или эмитент ждут нового допустимого действия.
Иногда это новый шаг аутентификации, иногда повтор в разрешённом окне,
иногда обработка кода Merchant Advice Code или уведомление клиента~\cite{visa-core-rules,mc-tpr}.

Если отказ требует аутентификации, клиента нужно вернуть в новый
подтверждённый шаг, а не пытаться задним числом \enquote{додавить} ту же
самую MIT. Европейское регулирование и схемные правила здесь
согласованы: при удалённом установлении MIT-мандата SCA применяется на
этапе создания мандата, а не после произвольной серии повторов~\cite{eba-mit-qa,mc-tpr}.

Если же речь не об аутентификации, а об отказе по регулярному списанию,
то действуют уже конкретные схемные процедуры. Visa для отклонённой MIT
с сохранённым реквизитом требует письменно уведомить держателя карты и
дать ему как минимум 7 календарных дней, чтобы заплатить другим
способом~\cite{visa-core-rules}. Mastercard для отказов по
регулярным платежам ограничивает повторную отправку (до 10~попыток
за 24~часа, до 35~попыток за 30~дней) и требует учитывать Merchant
Advice Code (MAC), который определяет допустимый интервал
и саму возможность повтора~\cite{mc-tpr}. Разница между лимитами
MC (35/30 дней) и Visa (15/30 дней) отражает разные методы подсчёта,
а не разную философию: оба лимита выставлены в зоне, где
совокупный уровень успеха повторов перестаёт расти, а нагрузка
на авторизационную инфраструктуру продолжает расти.

MAC определяет не только разрешение на повтор, но и конкретное
действие. Коды MAC~24--28 задают минимальное время ожидания перед
повтором: MAC~24 -- 1~час, MAC~25 -- 24~часа, MAC~26 -- 2~дня,
MAC~27 -- 6~дней, MAC~28 -- 10~дней. Эскалирующее окно устроено
не произвольно: код~24 рассчитан на временное ограничение баланса,
которое может сняться к следующему расчётному дню; код~28 -- на
ситуацию, когда держатель ждёт пополнения в следующем расчётном
цикле. Схема \enquote{обучает} мерчанта откладывать повтор до момента,
когда вероятность одобрения реально повысится. Отдельно стоит MAC~03
(\textit{Do not try again}), который полностью запрещает повторную
отправку по этим реквизитам, а MAC~01
(\textit{New account information available}) сигнализирует, что
реквизиты устарели и мерчант должен получить обновлённые данные
через ABU или от клиента~\cite{mc-tpr}. Игнорирование MAC
приводит к штрафному трафику: Mastercard отслеживает повторные
попытки после запрета и включает мерчанта в мониторинговую
программу.

Visa использует схожую логику через Response Code и Transaction
Advisor Code. Отказ с кодом, требующим аутентификации, возвращает
клиента в 3DS-цикл, а не в повторную MIT. Отказ с кодом,
допускающим повтор, ограничен 15~попытками за 30~дней для
одного мерчанта и набора реквизитов~\cite{visa-core-rules}.
Дальнейшие попытки Visa расценивает как нарушение правил.

Для биллинговой системы из этого вытекает конкретная задача:
каждый отказ должен проходить через классификатор, который
определяет, допустим ли повтор, через какое время и нужно ли
вернуть клиента в интерактивный сценарий. Без такого
классификатора биллинг либо повторяет слишком часто и попадает
под мониторинг схемы, либо сдаётся слишком рано и теряет
подписчиков.

\begin{figure}[tbp]
  \centering
  \includefiguresvg[width=.8\linewidth]{assets/figures/ch22-subscriptions/decline-decision-tree}
  \caption{Дерево решений при отказе MIT.}~\label{fig:decline-decision-tree}
\end{figure}

Отказ по MIT заставляет пересмотреть контекст платежа и выбрать
следующий шаг по правилам схемы.

\section{Обновление реквизитов}~\label{sec:subscriptions-account-updater}

Сервисы обновления реквизитов закрывают сценарий, когда клиент не
менял подписку, но эмитент перевыпустил карту. Сценарий знаком:
эмитент перевыпустил карту, старый PAN больше не подходит, а мерчант
продолжает пытаться списывать по старому набору данных.

Visa Account Updater (VAU) и Mastercard Automatic Billing Updater
(ABU) передают мерчанту обновлённые реквизиты от участвующих
эмитентов. Mastercard описывает ABU как программу, через которую
эквайер и мерчант получают обновлённые платёжные реквизиты, когда
карта истекла или была перевыпущена~\cite{mc-abu-privacy}; в самих
правилах Mastercard для регулярных платежей и операций с сохранёнными
реквизитами в ряде стран участие в ABU формализовано и сопровождается
исключениями по отдельным типам карт~\cite{mc-tpr}. Visa говорит об
этом шире как об управлении жизненным циклом реквизитов (credential
lifecycle management): обновление сохранённых реквизитов возможно в
реальном времени во время авторизации, причём речь идёт о PAN и о
сетевых токенах~\cite{visa-lifecycle-management}.

У этого подхода есть ограничения. Охват зависит от сети, рынка,
типа карты и модели участия. Не каждый карточный продукт попадает в
один и тот же контур обновления, а сами механизмы могут быть как
пакетными, так и встроенными прямо в авторизацию~\cite{mc-tpr,visa-lifecycle-management}.

Разница между пакетным и авторизационным режимом существенна
для архитектуры. Пакетный режим (batch) работает по расписанию:
мерчант или PSP отправляет файл с реквизитами, схема возвращает
обновлённые данные через несколько часов. Для подписок с месячным
циклом этого достаточно, если файл отправляется заблаговременно.
Авторизационный режим (real-time) встроен в сам запрос: если
реквизиты устарели, ответ содержит обновлённые данные, и мерчант
сохраняет их сразу. Visa описывает это как Real-Time Account
Updater~\cite{visa-lifecycle-management}; Mastercard -- как ABU
Real-Time Inquiry~\cite{mc-tpr}.

Покрытие зависит от участия эмитента. Не каждый банк подключён
к ABU или VAU, и не каждый тип карты попадает в контур
обновления. Visa рекомендует подключать все BIN, включая
reloadable prepaid, но нерезагружаемые prepaid-карты и продукты
малых эмитентов часто остаются за пределами сервиса. Обновление
реквизитов снижает технический отток, но не устраняет его:
для карт вне покрытия единственный путь -- уведомить клиента
и попросить ввести новые данные.

Сетевые токены меняют эту механику. Если мерчант хранит не PAN,
а токен схемы, привязка обновляется при перевыпуске карты
автоматически: токен остаётся прежним, underlying PAN обновляется
на стороне схемы~\cite{visa-lifecycle-management}. Это сдвигает
задачу обновления от мерчанта к схеме и делает batch-файлы
для токенизированных реквизитов ненужными.

\begin{figure}[tbp]
  \centering
  \includefiguresvg[width=.8\linewidth]{assets/figures/ch22-subscriptions/credential-update-flow}
  \caption{Обновление реквизитов: VAU/ABU и сетевые токены.}~\label{fig:credential-update-flow}
\end{figure}

Обновление реквизитов не заменяет дисциплину SCF и хранения
идентификатора исходной авторизации, но заметно снижает технический
отток и число отказов по старым картам.
Для зрелой подписочной системы это базовый элемент гигиены.

\section{Как выглядит устойчивый подписочный процесс}~\label{sec:subscriptions-anti-churn}

Устойчивый процесс начинается с первой CIT: она должна пройти,
корректно зафиксировав согласие. Затем система сохраняет идентификатор исходной
авторизации, токен или другой идентификатор реквизита, а клиент
получает ясное описание условий подписки:
периодичности, суммы и правил изменения.

Дальше каждый биллинговый цикл строится как MIT с корректными флагами,
своевременной проверкой актуальности реквизитов и понятной обработкой
отказов. Если карта перевыпущена, помогает обновление реквизитов. Если
эмитент просит вернуть клиента в цикл подтверждения, это надо сделать.
Если отказ носит временный характер, повторная попытка должна идти по
правилам схемы. Импровизация команды здесь особенно опасна, потому что
она размывает именно тот контекст, на котором держится вся подписка.

Мерчант с корректной NTID-цепочкой передаёт эмитенту другой сигнал:
не анонимный CNP-запрос, а продолжение известного мандата. Для
эмитента это снижает неопределённость: он применяет менее
строгую антифрод-политику. Практический результат: approval rate
на MIT с корректным NTID выше, чем на \enquote{сиротских} запросах без
истории цепочки.

Отдельная часть процесса касается коммуникации с клиентом. Предсказуемость
здесь ценнее красивого дизайна: клиент должен заранее знать,
когда будет следующая попытка списания и что именно ему нужно сделать,
если платёж не прошёл. Для подписки это часто важнее, чем красивая
форма уведомления, потому что понятное ожидание снижает и отток, и
число бессмысленных повторов.

\practicenote{%
Устойчивость подписки держится на четырёх вещах: корректная первая
CIT, идентификатор исходной авторизации, актуальные реквизиты и
правильная обработка отказов. Если один элемент выпадает, поток теряет контекст.

}
